\documentclass{jsarticle}
\usepackage{amssymb}
\usepackage{amsthm}
\usepackage{amsmath}
\newtheorem{theorem}{定理}
\newtheorem{proposition}[theorem]{命題}
\newtheorem{corollary}[theorem]{系}
\newtheorem{lemma}[theorem]{補題}
\newtheorem{defnition}[theorem]{定義}
\newtheorem{example}[theorem]{例}
\renewcommand{\proofname}{証明}
\newcommand{\card}{{\rm card}}
\title{始位相}
\author{電波通信}
\date{\today}
\begin{document}
\maketitle
この文書では始位相とその性質について紹介する。強く紹介したいのは定理\ref{main}である。この定理は筆者が位相空間論を初めて楽しく感じた定理である。始位相とはInitial topologyを訳したものだが逆位相とも呼ばれる。
\begin{defnition}[始位相]
$X$を集合$\{(Y_{\lambda},\mathfrak{O}_{\lambda})\}_{\lambda \in \Lambda}$を位相空間の族とし、各$\lambda$ 毎に写像
\[
f_{\lambda}:X\rightarrow Y_{\lambda}
\]が与えられているとするこの時$X$に写像族$F=\{f_{\lambda}\}_{\lambda \in \Lambda}$をすべて連続にするような最弱の位相が存在する。それを$F=\{f_{\lambda}\}_{\lambda \in \Lambda}$から誘導された始位相と呼ぶ。此処では記号$\mathfrak{I}(F)$で表すことにする。この位相は準開基$\mathfrak{M}=\{f_{\lambda}^{-1}(U)\ |\ U\in \mathfrak{O}_{\lambda},\ \lambda\in \Lambda \}$により生成されている
\end{defnition}$\mathfrak{I}$は$I$のフラクトゥールである。

\begin{example}[直積位相]
位相空間の族$\{(X_{\lambda},\mathfrak{O}_{\lambda})\}_{\lambda \in \Lambda}$の直積空間$\displaystyle \prod_{\lambda \in \Lambda}X_{\lambda}$の位相はその射影の族$\{\pi_{\lambda}\}_{\lambda \in \Lambda}$から定まる始位相である。
\end{example}

\begin{example}[相対位相]
相対位相は包含写像から定まる始位相である。
\end{example}

\begin{example}[シェルピンスキー空間と位相空間]シェルピンスキー空間とは集合$2=\{0,1\}$に
\[
\{\O,\{1\},\{0,1\}\}
\]
を開集合族として位相を定義した空間である。$(X,\mathfrak{O})$を任意の位相空間とするとこの位相は$\mathfrak{O}$によって添字付けられた開集合の特性関数の関数の族
\[
\chi_{U}:X\rightarrow \{0,1\}
\]によって定まる始位相である。より強く$(X,\mathfrak{O})$の開基の一つを$\mathfrak{B}$とするとこの空間の位相はシェルピンスキー空間への連続写像の族$\{\chi_{B}\}_{B\in \mathfrak{B}}$から定まる始位相である。証明は容易であろう。この事実は後で使う。
\end{example}

\begin{theorem}[始位相の普遍性]位相空間$Y$には$\{f_{\lambda}\}_{\lambda \in \Lambda}$から定まる始位相が入っているとしこの関数族の終域を$\{Z_{\lambda}\}_{\lambda \in \Lambda}$とするこの時位相空間$X$からの$Y$への写像$\phi$について
\[
\phi :X\rightarrow Yが連続\Leftrightarrow \forall \lambda\in \Lambda\ f_{\lambda}\circ \phi
: X\rightarrow Z_{\lambda}が連続
\]が成立する。
\end{theorem}
\begin{proof}
$[\Rightarrow]$の部分は明らかである。$[\Leftarrow]$を証明しよう。$X$の位相は準開基$\mathfrak{M}=\{f_{\lambda}^{-1}(U)\ |\ U\in \mathfrak{O}_{\lambda},\ \lambda\in \Lambda \}$により生成されていることを思い出そう。すると任意の$\lambda$と任意の$U\in \mathfrak{O}_{\lambda}$に対して
\[
\phi^{-1}(f_{\lambda}^{-1}(U))=(f_{\lambda}\circ \phi)^{-1}(U)
\]が開集合であることを示せば良いが、それは
\[
\forall \lambda\in \Lambda\ f_{\lambda}\circ \phi: X\rightarrow Z_{\lambda}が連続
\]であることから直ちに分かる。
\end{proof}


\begin{theorem}[始位相の埋め込み定理]\label{main}位相空間$(X,\mathfrak{O})$には$\{f_{\lambda}\}_{\lambda \in \Lambda}$から定まる始位相が入っているとしこの関数族の終域を$\{(Y_{\lambda},\mathfrak{O}_{\lambda})\}_{\lambda \in \Lambda}$とするこのとき関数
\[
\phi=\prod_{\lambda\in \Lambda}f_{\lambda}:X\rightarrow \prod_{\lambda\in \Lambda}Y_{\lambda \in \Lambda}
\]について$\mathfrak{I}(F)=\mathfrak{I}(\phi)$が成立する。
また$\phi$が単射ならばこれは埋め込みになっているつまり$(X,\mathfrak{O})$は$\prod_{\lambda\in \Lambda}Y_{\lambda \in \Lambda}$の部分空間$\phi(X)$と同相となる。
\end{theorem}
\begin{proof}${}$[前半]始位相の普遍性（直積位相ｎ普遍性を用いる）から$\phi$は$(X,\mathfrak{I}(F))$上$\phi$は連続なので最小性から
\[
\mathfrak{I}(\phi)\subset \mathfrak{I}(F)
\]そして$\pi_{\lambda}\circ \phi=F_{\lambda}$なので$f_{\lambda}$はに$(X,\mathfrak{I}(\phi)$上で連続なので始位相の最小性から
\[
\mathfrak{I}(F)\subset \mathfrak{I}(\phi)
\]
よって$\mathfrak{I}(F)=\mathfrak{I}(\phi)$が成り立つ。\\ ${}$
[後半]後半は$X$の位相が準開基
$\mathfrak{M}=\{f_{\lambda}^{-1}(U)\ |\ U\in \mathfrak{O}_{\lambda},\ \lambda\in \Lambda \}$
により生成されていることと今は$\phi$が単射であるから$\phi$の直像は交わりを保つ。よって$\phi(f_{\lambda}^{-1}(U))$が$\phi(X)$の開集合であることを示せばいい。此処で$f_{\lambda}=\pi _{\lambda}\circ \phi$を用いると
\[
f_{\lambda}^{-1}(U)=\phi^{-1}(\pi ^{-1}(U))
\]
と書くことができ、$\phi$の単射性を用いて
\[
\phi(f_{\lambda}^{-1}(U))=\phi(\phi^{-1}(\pi ^{-1}(U)))=\pi^{-1}(U)\cap \phi(X)
\]
を得て、直積位相と相対位相の定義から$\pi^{-1}(U)\cap \phi(X)$は$\phi(X)$の開集合である。これで証明が終わる。
\end{proof}
\begin{example}[$T_{0}$空間]$T_{0}$空間はシェルピンスキー空間の直積に埋め込み可能である。これは開基$\mathfrak{B}に対して$\\
\[
\prod_{U\in \mathfrak{B}}\chi_{U}が単射\Leftrightarrow (X,\mathfrak{O})がT_{0}
\]が成立するからである。
\end{example}
[$\displaystyle{\prod_{U\in \mathfrak{B}}\chi_{U}が単射\Leftrightarrow (X,\mathfrak{O})がT_{0}}$]の証明.\\ 
まず$T_0$とは\\ $(T_0)$\  異なる２点$x,y$について$x$の近傍で$y$を含まないようなものが存在するか、もしくは$y$の近傍で$x$を含まないようなものが存在する\\ 
であることを思い出そう。この定義の「近傍」を開基$\mathfrak{B}$の元としても良い。（開基なので）\\ 今、$n\neq y$と仮定すると$T_0$から($x$と$y$を適宜入れ替えて)$B\in \mathfrak{B}$が存在して$x\in B,\ y\not\in B$としても一般性を失わない。このとき
\[
\chi_B(x)=1,\ \chi_B(y)=0
\]なので
\[
\prod_{U\in \mathfrak{B}}\chi_{U}(x)\neq \prod_{U\in \mathfrak{B}}\chi_{U}(y)
\]よって
\[
x\neq y\Rightarrow \prod_{U\in \mathfrak{B}}\chi_{U}(x)\neq \prod_{U\in \mathfrak{B}}\chi_{U}(y)
\]なので$\prod_{U\in \mathfrak{B}}\chi_{U}$は単射。逆はこの証明を下から読めばわかる。
$\square$

\begin{defnition}[位相空間の重み]位相空間$X$の開基の最小濃度を位相空間$X$の重みと言い、$w(X)$で表す。つまり
\[
w(X)=\min\{\card(\mathfrak{B})\ |\ \mathfrak{B}はXの開基\}
\]
である。
\end{defnition}
この位相空間の重みを用いて、位相空間それ自体の濃度が評価出来る。
\begin{corollary}$T_0$空間に対して
\[
\text{\card$(X)$}\le2^{w(X)}
\]
\end{corollary}
\begin{proof}
$X$はシェルピンスキー空間の直積$\prod_{U\in \mathfrak{B}}\{0,1\}$に埋め込まれ、
\[
\text{card$(\prod_{U\in \mathfrak{B}}\{0,1\})$}=2^{\text{card$(\mathfrak{B})$}}
\]となるので、$\mathfrak{B}$を$\text{card$(\mathfrak{B})$}=w(X)$となるものとして選べば
\[
\text{card$(X)$}\le2^{w(X)}
\]が成り立つ。
\end{proof}
直ちに次のことがわかる。
\begin{corollary}第二可算公理を充たす$T_0$空間にたいして
\[
\card(X)\le \mathfrak{c}
\]

\end{corollary}

\begin{thebibliography}{9}
\bibitem{1}森田紀一,位相空間論,岩波全書331,岩波書店1981
\bibitem{2}M.G.Murdeshwar,general topology,Wiley Eastern Limited,1983
\end{thebibliography}










\end{document}